\subsection{Ideally Hausdorff modules}

DEFINITION 1. Let A be a commutative ring and \(\mathfrak{I}\) an ideal of A. An A-module M is called ideally Hausdorff with respect to \(\mathfrak{I}\) (or simply ideally Hausdorff if there is no ambiguity) if, for every finitely generated ideal \(\mathfrak{a}\) of A, the A-module \(\mathfrak{a} \otimes_{\mathbf{A}} \mathbf{M}\) is Hausdorff with the \(\mathfrak{I}\) -adic topology.

Putting \(\mathfrak{a} = \mathbf{A}\) in this definition, we have already seen that \(\mathbf{M}\) is necessarily Hausdorff with the 3-adic topology.

\textbf{Examples}\par

\begin{enumerate}
\def\labelenumi{(\arabic{enumi})}
\item
  If A is Noetherian and \(\mathfrak{I}\) is contained in the Jacobson radical of A (in other words if A is a Zariski ring with the \(\mathfrak{I}\) -adic topology), every finitely generated A-module is ideally Hausdorff ( \(\S\) 3, no. 3, Proposition 6).
\item
  Every direct sum of ideally Hausdorff modules is an ideally Hausdorff module, by virtue of the relations
\end{enumerate}

\[
\Im^ {n} \left(a \otimes_ {A} \bigoplus_ {\lambda \in L} M _ {\lambda}\right) = \Im^ {n} \bigoplus_ {\lambda \in L} (a \otimes_ {A} M _ {\lambda}) = \bigoplus_ {\lambda \in L} \Im^ {n} (a \otimes_ {A} M _ {\lambda}).
\]

\begin{enumerate}
\def\labelenumi{(\arabic{enumi})}
\setcounter{enumi}{2}
\tightlist
\item
  If an A-module M is flat and Hausdorff with the \(\Im\) -adic topology it is ideally Hausdorff, for \(a \otimes_{A} M\) is then identified with a submodule of M and the \(\Im\) -adic topology on \(a \otimes_{A} M\) is finer than the topology induced on \(a \otimes_{A} M\) by the \(\Im\) -adic topology on M, which is Hausdorff by hypothesis.
\end{enumerate}

